\chapter{Une image, une vidéo, vues par un ordinateur}
\label{chap:image-video}

\begin{quote}\itshape
Avant de parler d'intelligence artificielle, de vidéos cachées et de
modèles à deux milliards de réglages, mettons-nous d'accord sur une chose
toute simple : qu'est-ce qu'une image, pour une machine ? Ce chapitre pose
le décor, tranquillement, sans rien supposer.
\end{quote}

\begin{prerequis}
Rien du tout. Vraiment rien. Il faut savoir compter, savoir qu'une
multiplication sert à répéter une addition, et c'est tout. Chaque mot
nouveau sera expliqué au moment où il arrive.
\end{prerequis}

Vous regardez une photo de chat. Vous voyez un chat. C'est immédiat, c'est
gratuit, vous n'avez rien eu à faire. Un ordinateur, lui, ne voit pas de
chat. Il ne voit même pas d'image. Il voit une très longue liste de
nombres, et rien d'autre.

Tout ce livre part de là. Si nous voulons qu'une machine finisse par
comprendre une vidéo, il faut d'abord savoir sous quelle forme la vidéo
arrive chez elle. C'est exactement ce que nous allons faire dans ce
chapitre : ouvrir le capot et regarder les nombres.

\section{Un ordinateur ne voit pas des objets, il voit des nombres}

\begin{enclair}
Une image, pour un ordinateur, c'est une grille de petits carrés. Chaque
carré porte un nombre, et ce nombre dit une seule chose : à quel point ce
petit carré est clair ou sombre.
\end{enclair}

Pensez à une mosaïque, comme celles qu'on voit sur les murs des vieilles
maisons. De loin, vous voyez un dessin. De près, vous voyez des centaines
de petits carreaux collés les uns aux autres, chacun d'une seule couleur.
Une image numérique, c'est exactement cela : une mosaïque de petits
carrés, tous de la même taille, rangés en lignes et en colonnes.

\begin{definition}[Le pixel][pixel]
Un \emph{pixel} est un de ces petits carrés. C'est le plus petit morceau
d'une image : on ne peut pas le couper en deux. Le mot vient de l'anglais
\emph{picture element}, c'est-à-dire « élément d'image ». Chaque pixel
porte un nombre qui décrit sa couleur.
\end{definition}

Commençons par le cas le plus simple : une image en \emph{niveaux de
gris}, c'est-à-dire une image sans couleur, comme une vieille photo en
noir et blanc. Dans ce cas, un seul nombre suffit par pixel. On se met
d'accord une fois pour toutes sur ce que ce nombre veut dire :

\begin{itemize}
  \item le nombre \textbf{0} veut dire « ce carré est complètement noir » ;
  \item le nombre \textbf{255} veut dire « ce carré est complètement
        blanc » ;
  \item tous les nombres entre les deux sont des gris, de plus en plus
        clairs à mesure qu'on monte.
\end{itemize}

La figure~\ref{fig:image-grille} montre les deux faces de la même image : à
gauche ce que la machine a réellement en mémoire, à droite ce que nous
voyons apparaître à l'écran.

\begin{schema}{Les deux visages d'une même image. À gauche, les seize
nombres que l'ordinateur garde en mémoire. À droite, le dessin obtenu en
peignant chaque carré : petit nombre, carré sombre ; grand nombre, carré
clair. Il n'y a rien d'autre dans une image que ces nombres.}[image-grille]
  \begin{scope}[x=1cm, y=1cm]
    \node[font=\sffamily\small\bfseries, text=mbSlate] at (1.8,0.62)
         {Ce que la machine garde};
    \foreach \mbrow [count=\ri from 0] in
      {{240,235,210,245},{235,120,60,210},{215,60,20,120},{245,205,115,235}}{
      \foreach \mbval [count=\ci from 0] in \mbrow {
        \draw[draw=mbSlate!45, line width=0.5pt]
              (\ci*0.9,-\ri*0.9) rectangle ++(0.9,-0.9);
        \node[font=\small, text=mbInk]
              at (\ci*0.9+0.45,-\ri*0.9-0.45) {\mbval};
      }
    }
    \draw[fleche accent] (3.8,-1.8) -- (5.2,-1.8);
    \node[etiquette, above] at (4.5,-1.74) {on affiche};
    \begin{scope}[xshift=5.4cm]
      \foreach \mbrow [count=\ri from 0] in
        {{240,235,210,245},{235,120,60,210},{215,60,20,120},{245,205,115,235}}{
        \foreach \mbval [count=\ci from 0] in \mbrow {
          \pgfmathsetmacro{\mblum}{100-\mbval/2.55}
          \fill[black!\mblum] (\ci*0.9,-\ri*0.9) rectangle ++(0.9,-0.9);
        }
      }
      \draw[draw=mbSlate!45, line width=0.5pt] (0,0) rectangle (3.6,-3.6);
      \node[font=\sffamily\small\bfseries, text=mbSlate] at (1.8,0.62)
           {Ce que nous voyons};
    \end{scope}
  \end{scope}
\end{schema}

\begin{exemple}[Lire trois pixels]
Dans la grille de gauche de la figure~\ref{fig:image-grille}, le pixel tout en haut
à gauche vaut $240$ : il est presque blanc. Le pixel du milieu vaut $20$ :
il est presque noir. Un pixel qui vaudrait $128$, soit à peu près la
moitié de $255$, serait un gris moyen, ni clair ni foncé.
\end{exemple}

\begin{intuition}
Voyez chaque pixel comme un petit interrupteur à gradateur, du genre de
ceux qui règlent la lumière d'un salon. Tourné à fond dans un sens, la
lampe est éteinte : c'est $0$. Tourné à fond dans l'autre, la lampe éclaire
au maximum : c'est $255$. Une image, c'est un mur couvert de milliers de
ces petites lampes, chacune réglée à sa valeur.
\end{intuition}

\begin{piege}
Pourquoi $255$ et pas $100$, ou $1000$ ? C'est une simple convention, née
de la façon dont les ordinateurs rangent les nombres : avec huit cases de
mémoire, on peut écrire exactement $256$ valeurs différentes, de $0$ à
$255$. Retenez surtout le sens de lecture : \textbf{plus le nombre est
grand, plus le pixel est clair}. Beaucoup de débutants inversent, dessinent
des images en négatif, et cherchent l'erreur pendant une heure.
\end{piege}

Nous savons maintenant ce que contient une case de la mosaïque. La question
suivante vient toute seule : combien y a-t-il de cases ?

\section{La résolution : la finesse de la mosaïque}

\begin{definition}[La résolution][resolution]
La \emph{résolution} d'une image, c'est le nombre de pixels qu'elle
contient, donné sous la forme « largeur $\times$ hauteur ». Une image de
résolution $256 \times 256$ possède $256$ colonnes de pixels et $256$
lignes de pixels.
\end{definition}

Plus la mosaïque a de carreaux, plus le dessin est fin. Avec très peu de
carreaux, un rond ressemble à un escalier ; avec beaucoup, il redevient un
rond. La figure~\ref{fig:resolution} montre le même disque dessiné deux fois, la
première fois grossièrement, la seconde finement.

\begin{schema}{Le même disque, dessiné avec deux finesses différentes. À
gauche, la mosaïque est grossière et les bords deviennent des marches
d'escalier. À droite, les carreaux sont seize fois plus nombreux et l'\oe{}il
retrouve un vrai rond. Rien n'a changé dans le dessin : seule la finesse de
la grille a changé.}[resolution]
  \begin{scope}[x=1cm, y=1cm]
    % --- grille grossière : 6 x 6 ---
    \foreach \ri in {0,...,5}{
      \foreach \ci in {0,...,5}{
        \pgfmathsetmacro{\mbcx}{(\ci+0.5)/6-0.5}
        \pgfmathsetmacro{\mbcy}{(\ri+0.5)/6-0.5}
        \pgfmathsetmacro{\mbdd}{\mbcx*\mbcx+\mbcy*\mbcy}
        \pgfmathsetmacro{\mbin}{\mbdd<0.2100 ? 1 : 0}
        \ifdim\mbin pt>0.5pt
          \fill[mbBlue!65] (\ci*0.5,-\ri*0.5) rectangle ++(0.5,-0.5);
        \fi
        \draw[draw=mbSlate!30, line width=0.4pt]
              (\ci*0.5,-\ri*0.5) rectangle ++(0.5,-0.5);
      }
    }
    \node[font=\sffamily\small, text=mbSlate] at (1.5,-3.35)
         {$6 \times 6 = 36$ pixels};
    % --- grille fine : 24 x 24 ---
    \begin{scope}[xshift=5cm]
      \foreach \ri in {0,...,23}{
        \foreach \ci in {0,...,23}{
          \pgfmathsetmacro{\mbcx}{(\ci+0.5)/24-0.5}
          \pgfmathsetmacro{\mbcy}{(\ri+0.5)/24-0.5}
          \pgfmathsetmacro{\mbdd}{\mbcx*\mbcx+\mbcy*\mbcy}
          \pgfmathsetmacro{\mbin}{\mbdd<0.2100 ? 1 : 0}
          \ifdim\mbin pt>0.5pt
            \fill[mbBlue!65] (\ci*0.125,-\ri*0.125) rectangle ++(0.125,-0.125);
          \fi
        }
      }
      \draw[draw=mbSlate!30, line width=0.4pt] (0,0) rectangle (3,-3);
      \node[font=\sffamily\small, text=mbSlate] at (1.5,-3.35)
           {$24 \times 24 = 576$ pixels};
    \end{scope}
  \end{scope}
\end{schema}

Le nombre total de pixels s'obtient par une simple multiplication : on
multiplie la largeur par la hauteur, comme on calculerait le nombre de
carreaux d'un sol.

\begin{exemple}[Compter les pixels d'une vraie image]
Une image de résolution $256 \times 256$, qui est la taille utilisée
pendant la plus grande partie de l'entraînement de V-JEPA~2.1, contient
\[
  256 \times 256 = 65\,536 \text{ pixels.}
\]
Soit soixante-cinq mille petits carrés. Et ce n'est encore qu'une image en
gris : nous n'avons même pas parlé de couleur.
\end{exemple}

\begin{aretenir}
Deux images peuvent montrer exactement la même chose et ne pas avoir la
même résolution. Quand la résolution monte, le dessin gagne en finesse,
mais le nombre de nombres à traiter grimpe très vite : il grimpe comme le
\emph{carré} de la taille. Doubler la largeur et la hauteur, c'est
multiplier le travail par quatre. Cette remarque, qui paraît anodine ici,
expliquera plus tard pourquoi la dernière phase d'entraînement de
V-JEPA~2.1 est si courte.
\end{aretenir}

Nos images sont encore grises. Ajoutons la couleur : vous allez voir que
cela ne change presque rien à l'histoire.

\section{La couleur, c'est trois grilles au lieu d'une}

Un écran ne sait fabriquer que trois couleurs de lumière : du rouge, du
vert et du bleu. Toutes les autres couleurs sont obtenues en mélangeant ces
trois-là, en dosant chacune plus ou moins fort. Du rouge à fond plus du
vert à fond donne du jaune. Les trois à fond donnent du blanc. Les trois
éteintes donnent du noir.

Un pixel en couleur a donc besoin non pas d'un nombre, mais de trois : un
pour le rouge, un pour le vert, un pour le bleu. Chacun se lit comme avant,
de $0$ (éteint) à $255$ (à fond).

\begin{definition}[Le canal][canal]
Un \emph{canal} est une grille complète de nombres qui ne s'occupe que
d'une seule couleur de base. Une image en couleur possède trois canaux
superposés : le canal rouge, le canal vert et le canal bleu. On les appelle
souvent « les canaux R, V, B », ou « RGB » à l'anglaise.
\end{definition}

\begin{schema}{Une image en couleur, c'est trois mosaïques posées les unes
sur les autres. La première dit combien il y a de rouge dans chaque carré,
la deuxième combien de vert, la troisième combien de bleu. En regardant les
trois ensemble, l'\oe{}il reconstitue la couleur.}[canaux-rvb]
  \begin{scope}[x=1cm, y=1cm]
    \foreach \mbcol [count=\ri from 0] in {red, green, blue}{
      \begin{scope}[shift={(\ri*0.55,-\ri*0.5)}]
        \fill[\mbcol!18] (0,0) rectangle (2.4,2.4);
        \draw[\mbcol!70!black, line width=0.8pt] (0,0) rectangle (2.4,2.4);
        \foreach \mbk in {0.6,1.2,1.8}{
          \draw[\mbcol!35!white, line width=0.4pt] (\mbk,0) -- (\mbk,2.4);
          \draw[\mbcol!35!white, line width=0.4pt] (0,\mbk) -- (2.4,\mbk);
        }
      \end{scope}
    }
    \node[font=\sffamily\scriptsize\bfseries, text=red!60!black]
         at (2.75,2.15) {canal rouge};
    \node[font=\sffamily\scriptsize\bfseries, text=green!45!black]
         at (3.30,1.65) {canal vert};
    \node[font=\sffamily\scriptsize\bfseries, text=blue!60!black]
         at (3.85,1.15) {canal bleu};
    \draw[fleche accent] (4.9,0.7) -- (6.1,0.7);
    \node[etiquette, above] at (5.5,0.78) {on mélange};
    \begin{scope}[xshift=6.4cm, yshift=-0.5cm]
      \foreach \mbrow [count=\ri from 0] in
        {{blue!25,blue!25,blue!30,blue!25},
         {blue!25,orange!70,orange!80,blue!30},
         {blue!30,orange!80,orange!60,orange!70},
         {teal!35,teal!40,teal!35,teal!45}}{
        \foreach \mbval [count=\ci from 0] in \mbrow {
          \fill[\mbval] (\ci*0.6,{2.4-\ri*0.6}) rectangle ++(0.6,-0.6);
        }
      }
      \draw[draw=mbSlate!45, line width=0.6pt] (0,0) rectangle (2.4,2.4);
      \node[font=\sffamily\scriptsize\bfseries, text=mbSlate] at (1.2,-0.32)
           {l'image en couleur};
    \end{scope}
  \end{scope}
\end{schema}

\begin{exemple}[Quelques couleurs, en trois nombres]
Voici comment se lisent quelques couleurs courantes, dans l'ordre
rouge, vert, bleu :

\begin{center}
\begin{tabular}{lccc}
\toprule
\textbf{Couleur} & \textbf{Rouge} & \textbf{Vert} & \textbf{Bleu} \\
\midrule
Noir           & 0   & 0   & 0   \\
Blanc          & 255 & 255 & 255 \\
Rouge vif      & 255 & 0   & 0   \\
Jaune          & 255 & 255 & 0   \\
Gris moyen     & 128 & 128 & 128 \\
\bottomrule
\end{tabular}
\end{center}

Vous remarquez que le gris s'obtient en donnant la même valeur aux trois
canaux. C'est logique : aucune couleur ne tire la couverture à elle.
\end{exemple}

Comptons maintenant, pour de bon, le nombre de cases d'une image en
couleur.

\begin{exemple}[Le poids d'une image de $256 \times 256$]
Notre image possède $65\,536$ pixels, et chaque pixel réclame trois
nombres. Au total :
\[
  256 \times 256 \times 3 = 196\,608 \text{ nombres.}
\]
Presque deux cent mille nombres pour une seule petite image carrée, plus
petite que la moitié d'un écran de téléphone. Gardez ce chiffre en tête :
il va grossir très vite.
\end{exemple}

Nous avons beaucoup de nombres sur les bras. Il est temps de les ranger
proprement, et de leur donner un nom.

\section{Ranger les nombres : le tenseur et sa forme}

Pour parler de tous ces nombres sans se perdre, les informaticiens
utilisent un mot unique, qui a l'air savant mais qui ne cache rien de
compliqué : le \emph{tenseur}. Approchons-le par étapes, en montant une
marche à la fois.

\begin{schema}{La même idée, à quatre étages. Un nombre tout seul ; une
liste de nombres ; un tableau de nombres ; puis une pile de tableaux. Le
mot « tenseur » désigne tout cela d'un coup, quel que soit le nombre
d'étages.}[echelle-tenseur]
  \begin{scope}[x=1cm, y=1cm]
    % 1. un nombre
    \draw[draw=mbSlate!50, fill=mbBlue!10, line width=0.6pt]
          (0,-0.75) rectangle (0.5,-1.25);
    \node[font=\scriptsize, text=mbInk] at (0.25,-1) {7};
    \node[font=\sffamily\scriptsize, text=mbSlate, align=center]
          at (0.25,-2.9) {un nombre\\ (0 étage)};
    % 2. une liste
    \begin{scope}[xshift=2cm]
      \foreach \ci in {0,...,3}{
        \draw[draw=mbSlate!50, fill=mbBlue!10, line width=0.6pt]
              (\ci*0.5,-0.75) rectangle ++(0.5,-0.5);
      }
      \node[font=\sffamily\scriptsize, text=mbSlate, align=center]
            at (1,-2.9) {une liste\\ (1 étage)};
    \end{scope}
    % 3. un tableau
    \begin{scope}[xshift=5.2cm]
      \foreach \ri in {0,...,3}{
        \foreach \ci in {0,...,3}{
          \draw[draw=mbSlate!50, fill=mbBlue!10, line width=0.6pt]
                (\ci*0.5,-\ri*0.5) rectangle ++(0.5,-0.5);
        }
      }
      \node[font=\sffamily\scriptsize, text=mbSlate, align=center]
            at (1,-2.9) {un tableau\\ (2 étages)};
    \end{scope}
    % 4. une pile de tableaux
    \begin{scope}[xshift=8.6cm]
      \foreach \mbk in {2,1,0}{
        \begin{scope}[shift={(\mbk*0.28,-\mbk*0.28)}]
          \fill[mbAmber!\the\numexpr14+10*\mbk\relax] (0,0) rectangle (2,-2);
          \draw[draw=mbAmber!80!black, line width=0.6pt] (0,0) rectangle (2,-2);
          \foreach \mbj in {0.5,1,1.5}{
            \draw[draw=mbAmber!45!white, line width=0.4pt] (\mbj,0) -- (\mbj,-2);
            \draw[draw=mbAmber!45!white, line width=0.4pt] (0,-\mbj) -- (2,-\mbj);
          }
        \end{scope}
      }
      \node[font=\sffamily\scriptsize, text=mbSlate, align=center]
            at (1.3,-2.9) {une pile\\ (3 étages)};
    \end{scope}
  \end{scope}
\end{schema}

\begin{definition}[Le tenseur][tenseur]
Un \emph{tenseur} est un tableau de nombres qui peut avoir autant de
directions que l'on veut. Un nombre seul est un tenseur sans direction ;
une liste est un tenseur à une direction ; un tableau à lignes et colonnes
est un tenseur à deux directions ; une pile de tableaux est un tenseur à
trois directions, et ainsi de suite. Chaque direction s'appelle un
\emph{axe}.
\end{definition}

Pour décrire un tenseur, il suffit de dire combien de cases il compte le
long de chaque axe. Cette liste de tailles porte un nom, et c'est le mot le
plus utile de tout le chapitre.

\begin{notation}[La forme d'un tenseur][forme]
La \emph{forme} d'un tenseur est la liste du nombre de cases sur chaque
axe. On l'écrit entre parenthèses. Pour une image en couleur, on écrit
\[
  (\text{hauteur},\ \text{largeur},\ \text{canaux}),
\]
et notre image de tout à l'heure a donc la forme $(256, 256, 3)$.
\end{notation}

\begin{exemple}[Lire une forme]
Un tenseur de forme $(64, 384, 384, 3)$ se lit de gauche à droite comme une
phrase : « soixante-quatre choses, chacune haute de $384$ cases, large de
$384$ cases, et épaisse de $3$ canaux ». Nous verrons dans un instant que
ces soixante-quatre choses sont les images d'une vidéo.
\end{exemple}

\begin{piege}
L'ordre des axes n'est pas le même partout, et c'est la source d'erreur
numéro un des débutants. Deux conventions se partagent le monde :

\begin{itemize}
  \item $(\text{hauteur},\ \text{largeur},\ \text{canaux})$ — l'ordre des
        outils d'affichage et de photo ;
  \item $(\text{canaux},\ \text{hauteur},\ \text{largeur})$ — l'ordre des
        bibliothèques d'apprentissage comme PyTorch.
\end{itemize}

Les nombres sont exactement les mêmes, seule leur façon d'être rangés
change. Mais si vous vous trompez, votre programme croira qu'une image de
$256 \times 256$ en couleur est une image de $3$ pixels de haut sur $256$ de
large avec $256$ canaux, et vous obtiendrez une bouillie. Prenez
l'habitude, dès aujourd'hui, d'écrire la forme attendue en commentaire
au-dessus de chaque ligne de code qui manipule une image.
\end{piege}

Nous savons ranger une image. Passons à ce qui nous intéresse vraiment dans
ce livre : la vidéo.

\section{Une vidéo, c'est une pile d'images}

Une vidéo n'est rien d'autre qu'une suite d'images fixes, montrées les unes
après les autres, très vite. Si elles défilent assez vite, notre \oe{}il ne
voit plus des images séparées : il voit du mouvement. C'est le principe du
petit carnet à feuilleter, celui où l'on dessine un bonhomme dans le coin
de chaque page et où l'on fait défiler les pages avec le pouce.

\begin{definition}[L'image d'une vidéo][frame]
Chacune des images fixes qui composent une vidéo s'appelle une
\emph{image}, ou une \emph{trame} ; en anglais on dit \emph{frame}, et ce
mot revient tellement souvent dans les articles et dans le code qu'il faut
le connaître. Une vidéo est donc une pile d'images empilées non pas dans
l'espace, mais dans le temps.
\end{definition}

\begin{definition}[Les images par seconde][fps]
Le nombre d'\emph{images par seconde} dit combien d'images sont montrées
pendant une seconde. On l'abrège souvent en \emph{fps}, de l'anglais
\emph{frames per second}. Un film de cinéma tourne à $24$ images par
seconde ; les vidéos de téléphone, souvent à $30$ ou $60$.
\end{definition}

\begin{schema}{Une vidéo vue comme une pile de feuilles. Chaque feuille est
une image complète, avec ses trois canaux de couleur. La flèche du bas est
le temps : elle avance d'un cran à chaque nouvelle feuille. Ici, quatre
images par seconde, donc un quart de seconde entre deux
feuilles.}[pile-video]
  \begin{scope}[x=1cm, y=1cm]
    \foreach \mbk in {4,3,2,1,0}{
      \begin{scope}[shift={(\mbk*1.35,\mbk*0.42)}]
        \fill[mbBlue!\the\numexpr12+8*\mbk\relax] (0,0) rectangle (1.9,1.35);
        \draw[draw=mbBlue!75, line width=0.7pt] (0,0) rectangle (1.9,1.35);
        \foreach \mbj in {0.475,0.95,1.425}{
          \draw[draw=mbBlue!30!white, line width=0.35pt] (\mbj,0) -- (\mbj,1.35);
        }
        \foreach \mbj in {0.45,0.9}{
          \draw[draw=mbBlue!30!white, line width=0.35pt] (0,\mbj) -- (1.9,\mbj);
        }
      \end{scope}
    }
    \draw[-{Stealth[round,length=2.6mm,width=2mm]}, draw=mbAmber!85!black,
          line width=1pt] (0.2,-0.55) -- (7.6,-0.55);
    \foreach \mbk in {0,1,2,3,4}{
      \draw[draw=mbAmber!85!black, line width=0.7pt]
            (\mbk*1.35+0.95,-0.45) -- (\mbk*1.35+0.95,-0.65);
    }
    \node[font=\sffamily\scriptsize, text=mbSlate] at (0.95,-0.95) {$0$~s};
    \node[font=\sffamily\scriptsize, text=mbSlate] at (2.30,-0.95) {$0{,}25$~s};
    \node[font=\sffamily\scriptsize, text=mbSlate] at (3.65,-0.95) {$0{,}50$~s};
    \node[font=\sffamily\scriptsize, text=mbSlate] at (5.00,-0.95) {$0{,}75$~s};
    \node[font=\sffamily\scriptsize, text=mbSlate] at (6.35,-0.95) {$1$~s};
    \node[font=\sffamily\small\bfseries, text=mbAmber!80!black]
          at (7.9,-0.95) {temps};
    \node[annotation bulle] (mbune) at (-0.9,2.6)
         {une image\\ $(H, W, 3)$};
    \draw[trait annotation] (mbune) to[bend right=15] (0.55,1.0);
  \end{scope}
\end{schema}

Puisque la vidéo est une pile d'images, sa forme s'obtient en ajoutant un
axe devant celle d'une image : l'axe du temps.

\begin{notation}[La forme d'une vidéo][forme-video]
Un extrait de vidéo s'écrit sous la forme
\[
  (\text{temps},\ \text{hauteur},\ \text{largeur},\ \text{canaux}),
\]
où le premier nombre est le nombre d'images de l'extrait. On appelle
souvent un tel extrait un \emph{clip} : c'est un petit morceau découpé dans
une vidéo plus longue, exactement comme on découpe un morceau dans une
bobine de film.
\end{notation}

La durée d'un clip se calcule sans effort. Si vous avez un certain nombre
d'images, et que vous savez combien d'images passent en une seconde, alors
la division vous donne les secondes :
\[
  \text{durée en secondes} \;=\; \frac{\text{nombre d'images}}
                                      {\text{images par seconde}}.
\]

\begin{exemple}[La durée d'un clip d'entraînement]
Pendant la plus grande partie de son entraînement, V-JEPA~2.1 regarde des
clips de $16$ images prises à raison de $4$ images par seconde. Ces clips
durent donc
\[
  \frac{16}{4} = 4 \text{ secondes.}
\]
Quatre secondes, c'est court, mais c'est bien assez pour voir un geste
commencer et se terminer : une main qui saisit une tasse, une porte qui
s'ouvre. Et c'est justement ce que le modèle doit apprendre à comprendre.
\end{exemple}

\begin{surlesdonnees}
L'article de V-JEPA~2.1 utilise deux réglages, l'un après l'autre. Pendant
la longue phase d'entraînement principale, les clips font $16$ images en
$256 \times 256$. Pendant la courte phase finale, appelée phase de
refroidissement, on monte à $64$ images en $384 \times 384$. Toujours à
quatre images par seconde, ces clips-là durent
\[
  \frac{64}{4} = 16 \text{ secondes.}
\]
Le modèle regarde donc, à la fin, des extraits quatre fois plus longs et
beaucoup plus fins. Nous verrons au chapitre consacré au refroidissement
pourquoi cet ordre — d'abord grossier et court, ensuite fin et long — n'est
pas un caprice.
\end{surlesdonnees}

Comptons enfin les nombres contenus dans un clip, en multipliant tout
simplement les quatre tailles entre elles.

\begin{exemple}[Le nombre de cases d'un clip]
Pour un clip de $16$ images en $256 \times 256$ couleur :
\[
  16 \times 256 \times 256 \times 3 = 3\,145\,728 \text{ nombres.}
\]
Plus de trois millions. Et pour un clip de la phase de refroidissement,
$64$ images en $384 \times 384$ :
\[
  64 \times 384 \times 384 \times 3 = 28\,311\,552 \text{ nombres,}
\]
soit environ vingt-huit millions, c'est-à-dire neuf fois plus que le
précédent.
\end{exemple}

\begin{aretenir}
Une image est un tenseur à trois axes, une vidéo un tenseur à quatre axes.
Rien de plus mystérieux : on a seulement ajouté le temps devant. Toutes les
formes que vous croiserez dans ce livre se lisent avec cette seule règle.
\end{aretenir}

Trois millions de nombres, cela reste abstrait. Regardons ce que cela
représente vraiment dans la mémoire d'une machine.

\section{Ce que tout cela pèse vraiment}

Un ordinateur range ses nombres dans des cases de mémoire. L'unité de
mesure de cette mémoire est l'\emph{octet}.

\begin{definition}[L'octet][octet]
Un \emph{octet} est un petit paquet de mémoire, celui qui suffit tout juste
à écrire un nombre entier compris entre $0$ et $255$. C'est aussi l'unité
dans laquelle on mesure la taille des fichiers. Mille octets font environ un
kilo-octet, un million d'octets font environ un méga-octet (noté Mo), un
milliard d'octets font environ un giga-octet (noté Go).
\end{definition}

Nos pixels vont de $0$ à $255$ : ils tiennent donc chacun dans un seul
octet. Mais dès qu'un modèle se met à calculer, il a besoin de nombres à
virgule — parce qu'une moyenne, une multiplication ou une petite correction
tombent rarement sur un compte rond. Ces nombres à virgule s'appellent des
\emph{nombres flottants}, et le format le plus courant en occupe quatre
octets chacun, soit quatre fois plus.

\begin{exemple}[Du fichier à la mémoire de calcul]
Notre clip de $3\,145\,728$ nombres pèse :

\begin{itemize}
  \item $3\,145\,728$ octets, soit environ \textbf{3 Mo}, tant que les
        pixels restent des entiers de $0$ à $255$ ;
  \item $3\,145\,728 \times 4 = 12\,582\,912$ octets, soit environ
        \textbf{12,6 Mo}, une fois convertis en nombres à virgule pour le
        calcul.
\end{itemize}
\end{exemple}

Or un modèle ne regarde jamais un seul clip à la fois : il en regarde tout
un paquet d'un coup, pour aller plus vite. Ce paquet s'appelle un
\emph{lot}, et sa taille est un réglage du programme d'entraînement.

\begin{exemple}[Le poids d'un lot entier]
V-JEPA~2.1 traite $128$ clips à la fois. Pour la phase principale, cela
donne
\[
  128 \times 12\,582\,912 \approx 1\,610\,000\,000 \text{ octets,}
\]
soit environ \textbf{1,6 Go} — rien que pour les vidéos d'entrée, avant
même le moindre calcul. Pour la phase de refroidissement, où chaque clip
pèse neuf fois plus, on arrive à environ \textbf{14,5 Go}. Voilà, en une
ligne de calcul, pourquoi cette phase-là ne dure que très peu de temps :
elle coûte cher.
\end{exemple}

\begin{enia}
Ce chapitre a l'air de ne parler que de pixels, mais il pose déjà les trois
problèmes que tout le reste du livre s'emploie à résoudre.

\begin{itemize}
  \item \textbf{Il y a beaucoup trop de nombres.} Aucun modèle ne peut
        traiter trois millions de valeurs une par une. Il faudra donc
        d'abord les regrouper en petits carrés de $16 \times 16$ pixels :
        c'est ce que nous ferons au chapitre~4, consacré au transformeur de
        vision.
  \item \textbf{La finesse coûte cher.} Chaque fois que la résolution
        monte, la facture explose. C'est pour cela que l'entraînement
        commence en basse résolution et ne monte qu'à la toute fin.
  \item \textbf{Lire une vidéo prend du temps.} Une vidéo est stockée
        compressée ; il faut la décompresser image par image avant de
        pouvoir en faire quoi que ce soit. En production, ce travail-là
        devient souvent plus lent que le modèle lui-même. Nous y
        reviendrons longuement dans la partie consacrée à la mise en
        service.
\end{itemize}
\end{enia}

\begin{plusloin}
Dans le même temps où il regarde $128$ clips, V-JEPA~2.1 regarde aussi
$2\,304$ images fixes. En reprenant nos calculs, ces images pèsent
$2\,304 \times 196\,608 \times 4 \approx 1,8$~Go, c'est-à-dire à peu près
autant que les $1,6$~Go de vidéo. Ces deux nombres, $128$ et $2\,304$,
n'ont donc pas été tirés au hasard : ils ont été choisis ensemble pour que
les deux moitiés du travail coûtent à peu près pareil. Nous y reviendrons
au chapitre sur les données.
\end{plusloin}

Assez de calculs sur le papier. Vérifions tout cela sur une vraie machine.

\section{Passons au code}

Nous allons manipuler nos images et nos clips avec Python, le langage
utilisé par presque tous les travaux d'intelligence artificielle, et avec
\co{python}{numpy}, la bibliothèque qui sait manipuler les tenseurs. Vous
n'avez pas besoin de comprendre chaque ligne dès maintenant : lisez d'abord
les commentaires, et comparez les nombres affichés avec ceux que nous
venons de calculer à la main.

\begin{code}{python}[Fabriquer une image et lire sa forme]
import numpy as np

# Une image en couleur de 256 x 256 pixels, entierement noire.
# La forme se lit : (hauteur, largeur, canaux).
image = np.zeros((256, 256, 3), dtype=np.uint8)

print(image.shape)   # -> (256, 256, 3)
print(image.size)    # -> 196608   le nombre total de cases
print(image.dtype)   # -> uint8    un entier de 0 a 255
print(image.nbytes)  # -> 196608   un octet par case

# Peignons le pixel de la ligne 10, colonne 20, en rouge vif.
image[10, 20] = [255, 0, 0]
print(image[10, 20])  # -> [255   0   0]
\end{code}

Le mot \co{python}{uint8} se lit « entier non signé sur huit bits » : c'est
exactement le type de nombre qui va de $0$ à $255$, celui dont nous avons
parlé. Passons maintenant à la vidéo.

\begin{code}{python}[Un clip, sa forme, sa durée et son poids]
import numpy as np

# Un clip de 16 images de 256 x 256 en couleur.
# La forme se lit : (temps, hauteur, largeur, canaux).
clip = np.zeros((16, 256, 256, 3), dtype=np.uint8)

print(clip.shape)   # -> (16, 256, 256, 3)
print(clip.size)    # -> 3145728   les trois millions annonces

images_par_seconde = 4.0
duree = clip.shape[0] / images_par_seconde
print(duree, "secondes")            # -> 4.0 secondes

# Le meme clip converti en nombres a virgule, ramenes entre 0 et 1.
clip_flottant = clip.astype(np.float32) / 255.0
print(clip_flottant.nbytes / 1e6, "Mo")   # -> 12.582912 Mo
\end{code}

Enfin, voici le changement d'ordre des axes dont nous parlions plus haut.
C'est une opération que vous ferez des centaines de fois, alors autant la
voir tout de suite.

\begin{code}{python}[Passer de l'ordre « affichage » à l'ordre « PyTorch »]
import numpy as np
import torch

clip = np.zeros((16, 256, 256, 3), dtype=np.uint8)   # (T, H, W, C)

x = torch.from_numpy(clip).float() / 255.0
x = x.permute(3, 0, 1, 2)                            # (C, T, H, W)

print(x.shape)   # -> torch.Size([3, 16, 256, 256])

# Aucun nombre n'a ete perdu : seul l'ordre des axes a change.
print(x.numel())  # -> 3145728
\end{code}

\begin{piege}
Remarquez la dernière ligne : le nombre de cases n'a pas bougé. Changer
l'ordre des axes ne crée et ne détruit aucune donnée, cela réorganise
seulement la façon de les parcourir. Si jamais ce nombre change après une
telle opération, c'est que vous avez commis une erreur ailleurs — et cette
vérification en une ligne vous fera gagner des heures.
\end{piege}

\begin{aretenir}
Une image est une grille de nombres ; la couleur ajoute une troisième
dimension ; la vidéo en ajoute une quatrième, le temps. Tout le reste de ce
livre travaille sur ces tenseurs-là, et sur rien d'autre. Quand vous serez
perdu dans un chapitre plus difficile, revenez à cette question toute
simple : « quelle est la forme de ce que je manipule ? » Elle remet presque
toujours les idées en place.
\end{aretenir}

Nous savons désormais sous quelle forme une image et une vidéo arrivent
dans une machine. Mais nous n'avons encore rien appris à personne : nos
tenseurs restent inertes. La question suivante est donc la bonne :
qu'est-ce que cela veut dire, au juste, qu'une machine « apprend » ? C'est
tout le sujet du chapitre suivant.

\section*{Exercices}
\addcontentsline{toc}{section}{Exercices}

\begin{exercice}[Compter les cases]
Combien de nombres contient une image en couleur de $512 \times 512$
pixels ? C'est la taille des images utilisées pendant la phase de
refroidissement de V-JEPA~2.1.
\end{exercice}

\begin{solution}
On multiplie la hauteur, la largeur et le nombre de canaux :
$512 \times 512 \times 3 = 262\,144 \times 3 = 786\,432$ nombres. C'est
quatre fois plus qu'une image de $256 \times 256$, ce qui est normal
puisque nous avons doublé les deux côtés.
\end{solution}

\begin{exercice}[La durée d'un clip]
Un clip contient $32$ images, prises à raison de $8$ images par seconde.
Combien de temps dure-t-il ? Et combien de temps durerait-il si les mêmes
$32$ images étaient prises à $4$ images par seconde ?
\end{exercice}

\begin{solution}
Dans le premier cas, $32 / 8 = 4$ secondes. Dans le second,
$32 / 4 = 8$ secondes. Le nombre d'images est le même, mais elles sont
prélevées plus loin les unes des autres : le clip couvre donc une tranche
de temps deux fois plus longue. C'est un réglage important, car il décide
de ce que le modèle a le droit de voir.
\end{solution}

\begin{exercice}[Le poids d'un lot]
Un lot contient $64$ clips de $16$ images en $256 \times 256$ couleur,
convertis en nombres à virgule de quatre octets. Combien de méga-octets
occupe ce lot ?
\end{exercice}

\begin{solution}
Un clip contient $16 \times 256 \times 256 \times 3 = 3\,145\,728$ nombres,
donc $3\,145\,728 \times 4 = 12\,582\,912$ octets, soit environ $12{,}6$~Mo.
Pour $64$ clips : $64 \times 12{,}6 \approx 806$~Mo, c'est-à-dire un peu
plus de huit cents méga-octets. Ce lot fait exactement la moitié de celui
de l'article, qui en contient $128$, et l'on retrouve bien les $1{,}6$~Go
calculés plus haut en doublant ce résultat.
\end{solution}

\begin{exercice}[Lire une forme]
Un collègue vous tend un tenseur de forme $(3, 64, 384, 384)$ et vous dit
que c'est une vidéo. Que représente chaque nombre ? Quelle est la
convention utilisée, et quelle serait la forme du même tenseur dans l'autre
convention ?
\end{exercice}

\begin{solution}
Le $3$ ne peut être que le nombre de canaux de couleur : aucune vidéo ne
fait trois pixels de haut. Il s'agit donc de l'ordre
$(\text{canaux},\ \text{temps},\ \text{hauteur},\ \text{largeur})$, celui de
PyTorch : une vidéo en couleur de $64$ images en $384 \times 384$. Dans la
convention d'affichage, le même tenseur s'écrirait
$(64, 384, 384, 3)$. Les nombres qu'il contient sont exactement les mêmes ;
seul leur rangement change.
\end{solution}
